\section{Tres dominios de aplicación representativos}

\subsection{Coding Agent: el repositorio de código como entorno ejecutable}

Las tareas de programación cuentan con retroalimentación natural: el compilador, las pruebas, el análisis estático y los registros de ejecución pueden funcionar como verifier. El valor de un coding agent no es solo completar código, sino mantener un estado de trabajo en constante cambio: buscar en el repositorio, editar múltiples archivos, ejecutar pruebas, leer los fallos, ubicar la causa raíz e iterar.

\begin{figure}[H]
\centering
\includegraphics[width=0.94\textwidth]{figures/coding-agents.jpg}
\caption{Un coding agent opera simultáneamente el shell, el navegador, el editor y el estado del plan.\protect\footnotemark}
\end{figure}
\footnotetext{Intervalo de la explicación en video: 00:52:44--00:54:34.}

\begin{knowledgebox}{Por qué el código es el terreno de prueba ideal para métodos de automejora}
Las tareas de código combinan generación abierta con verificación automática relativamente barata: el espacio de programas candidatos es enorme, pero las pruebas unitarias, la verificación de tipos y los resultados de ejecución permiten filtrar candidatos con rapidez. Esto convierte al código en un punto de encuentro de alto valor para repeated sampling, search, execution feedback y RL.
\end{knowledgebox}

\subsection{Customer Support Agent: la retroalimentación proviene del estado del negocio}

Un agente de atención al cliente necesita acceso a pedidos, tickets, reembolsos y bases de conocimiento, y debe ejecutar acciones reales del negocio. El criterio de éxito normalmente no es que "la respuesta suene razonable", sino si el problema se resolvió, si el estado se actualizó correctamente y si el usuario necesita volver a contactar al servicio.

\begin{figure}[H]
\centering
\includegraphics[width=0.93\textwidth]{figures/customer-support-agents.jpg}
\caption{Un agente de atención al cliente conecta la recuperación de información, el juicio y las acciones del negocio.\protect\footnotemark}
\end{figure}
\footnotetext{Intervalo de la explicación en video: 00:54:34--00:55:43.}

La dificultad de este tipo de sistemas radica en los permisos y las acciones irreversibles: el modelo puede sugerir un reembolso, pero antes de ejecutarlo realmente suele requerirse verificación de reglas, límites de monto o confirmación humana.

\subsection{Research Agent: dividir el proceso de investigación en etapas verificables}

Un agente de investigación abarca la generación de ideas, la búsqueda bibliográfica, el diseño experimental, la ejecución de código, el análisis de resultados y la redacción del artículo. El verifier varía mucho entre etapas: en la búsqueda se puede comprobar si las citas existen, en el código se pueden correr pruebas, pero las conclusiones experimentales requieren rigor estadístico y juicio del dominio.

\begin{figure}[H]
\centering
\includegraphics[width=0.96\textwidth]{figures/research-agents.jpg}
\caption{El flujo de ideas, iteración experimental y redacción del artículo de un research agent.\protect\footnotemark}
\end{figure}
\footnotetext{Intervalo de la explicación en video: 00:56:00--00:59:03.}

\begin{warningbox}{"Completar automáticamente el informe" no equivale a completar la investigación}
Un research agent es especialmente propenso a producir resultados que parecen completos en la superficie pero con evidencia débil en cuanto a la veracidad de las citas, la reproducibilidad experimental y la explicación causal. Un sistema de alta calidad debe incorporar en su flujo de trabajo la verificación de fuentes, los registros de experimentos, las pruebas estadísticas y la búsqueda de contraejemplos, en lugar de optimizar únicamente la fluidez del texto final.
\end{warningbox}

\subsection{Resumen de la sección}

El punto en común de los tres tipos de aplicaciones es que el modelo debe salir del espacio puramente textual e interactuar con un estado externo cambiante. Sus diferencias provienen principalmente del costo y la fiabilidad del verifier: la retroalimentación de la ejecución de código es la más barata, la señal de éxito del negocio llega con más retraso, y la corrección de la investigación es la más difícil de automatizar.
